\documentclass[10pt,a4paper]{amsart}
\usepackage[margin=19mm]{geometry}
\usepackage{amsmath,amssymb,mathtools}
\usepackage[scr=rsfs,cal=euler]{mathalfa}
\usepackage{xfrac}
\usepackage[unicode,hidelinks,bookmarks=false]{hyperref}
\newcommand{\ie}{i.e.\ }
\newcommand{\cf}{cf.\ }
\newcommand{\resp}{respectively\ }
\newcommand{\Subsection}{\subsection}
\providecommand{\nobreakdash}{\nobreak\hbox{-}}
\setlength{\emergencystretch}{2em}
\multlinegap0pt
\usepackage{amssymb}
\newtheorem{theorem}{Theorem}
\title{Variable Calder\'on-Hardy spaces on the Heisenberg group}
\author{Pablo Rocha}
\date{Source: arXiv:2505.15760v2 (CC BY 4.0)}
\begin{document}
\maketitle

\noindent\emph{Source and licence.} Variable Calder\'on-Hardy spaces on the Heisenberg group. Version arXiv:2505.15760v2 (CC BY 4.0), distributed by arXiv under the Creative Commons Attribution 4.0 International licence, http://creativecommons.org/licenses/by/4.0/, as recorded in the arXiv licence metadata for this exact version and shown on the abstract page https://arxiv.org/abs/2505.15760v2. Full licence notice: \texttt{licenses/arxiv-2505.15760-CC-BY-4.0.txt}.\par\smallskip

\noindent\textit{Editorial scope.} Counted unit: the article's theorem thm (source0.tex lines 795--798). The article's complete original proof of it is delivered with it (source0.tex lines 800--817, one \texttt{proof} environment of 18 lines). No mathematics is written, repaired or re-derived: the statement and the proof are the article's own lines, in the article's own notation, and no retained line is substituted or re-flowed.  No further source-local context is needed: every cross-reference of every retained line resolves inside the two delivered spans.  Nothing was omitted from any retained span, so the package carries no omission record.  No retained line contains a citation, so the wrapper carries no bibliography and no bibliography entry.  No retained line contains a figure environment, an image-inclusion command or drawing code, so no image is delivered and the package declares no asset dependency.  Editorial formatting only, in addition: an AMS \texttt{amsart} preamble that restates the article's own declarations and macros used by the retained lines, namely the environments \texttt{theorem} on separate counters, together with the standard packages those lines need.  The retained environments print in this document as Theorem 1 in order of appearance, whereas the article prints the same items with the numbering of its own full document.  The complete native source of the article as distributed in the arXiv e-print for this exact version is bundled unchanged as \texttt{sources/178536/source0.tex}.
\begin{theorem} \label{second thm}
If $p(\cdot)$ is an exponent function on $\mathbb{H}^n$ such that $p_{+} \leq Q (2 + Q/q)^{-1}$, then 
$\mathcal{H}^{p(\cdot)}_{q, \, 2}(\mathbb{H}^{n}) = \{ 0 \}.$
\end{theorem}

\begin{proof} Let $G \in \mathcal{H}^{p(\cdot)}_{q, \, 2}(\mathbb{H}^{n})$ and assume $G \neq 0$. Then there exists 
$g \in G$ that is not a polynomial of homogeneous degree less or equal to $1$. It is easy to check that there exist a positive constant $c$ and a $\rho-$ball $B = B(e,r)$ with $r>1$ such that
\[
\int_{B} |g(w) - P(w)|^{q} \, dw \geq c > 0,
\]
for every $P \in \mathcal{P}_{1}$.

Let $z$ be a point such that $\rho(z) > r$ and let $\delta = 2 \rho(z)$. Then $B(e, r) \subset B(z, \delta)$. If $h \in G$, then 
$h = g - P$ for some $P \in \mathcal{P}_{1}$ and
\[
\delta^{-2}|h|_{q, B(z, \delta)} \geq c \rho(z)^{-2-Q/q}.
\]
So $N_{q,2}(G; \, z) \geq c \, \rho(z)^{-2-Q/q}$, for $\rho(z) > r$. Since $p_{+} \leq Q(2+Q/q)^{-1}$, we have that
\[
\kappa_{p(\cdot)} \left( N_{q,2m}(G; \cdot) \right) \geq c \, \int_{\rho(z) > r} \rho(z)^{-(2+Q/q)p_{+}} \, dz = \infty,
\]
which gives a contradiction. Thus $\mathcal{H}^{p}_{q, 2}(\mathbb{H}^{n}) = \{0\}$, if $p_{+} \leq Q(2+Q/q)^{-1}$.
\end{proof}

\end{document}
